\documentclass[a4paper,11pt]{scrartcl}

\usepackage[T1]{fontenc}
\usepackage[utf8]{inputenc}
\usepackage[ngerman]{babel}
\usepackage{lmodern}
\usepackage{amsmath}
\usepackage{amssymb}
\usepackage[a4paper,margin=2.5cm]{geometry}
\usepackage{enumitem}
\usepackage{booktabs}
\usepackage[hidelinks]{hyperref}
\usepackage{catppuccin-dark}

\newcounter{aufgabe}
\newcommand{\Aufgabe}[2]{%
  \refstepcounter{aufgabe}%
  \subsection*{Aufgabe~\theaufgabe: #1 \hfill (#2~Punkte)}%
  \label{auf:\theaufgabe}}

\title{Übungsblatt 4: Newton Methods (Newton-Verfahren)}
\subtitle{Begleitmaterial zur Vorlesung \glqq Numerical Methods\grqq{} (Prof.\ Dr.-Ing.\ Mark Schutera)}
\author{Newton-Raphson, Taylor-Entwicklung und Sekantenverfahren\\(Skript S.~35--46)}
\date{\today}

\begin{document}

\maketitle

\noindent\textbf{Gesamt: 50 Punkte.}\quad
Hilfsmittel: Taschenrechner. Runden Sie bei den Handrechnungen auf sieben Nachkommastellen
(Endergebnisse mit vier Nachkommastellen genügen, sofern nicht anders angegeben).
Die Aufgaben steigen im Schwierigkeitsgrad an: Aufgaben 1--2 prüfen das Begriffs- und
Herleitungsverständnis, Aufgaben 3--6 sind Rechenaufgaben analog zum Skript (mit eigenen
Beispielen), Aufgaben 7--8 sind Transferaufgaben. Die Gleichungsnummern (20)--(34) beziehen
sich auf das Skript.

\bigskip
\hrule

\section*{Aufgaben}

\Aufgabe{Lokale Information statt blinder Suche (local information vs.\ grid search)}{4}

Das Skript motiviert Newtons Methode als fundamentale Abkehr von der blinden
Brute-Force-Gittersuche (grid search).

\begin{enumerate}[label=(\alph*)]
  \item Was ist der Hauptvorteil von Newtons Methode gegenüber Grid Search? Formulieren Sie
  die \emph{Schlüsselidee} (key insight) des Skripts in eigenen Worten: Welche \emph{zwei}
  Fragen stellt Newtons Methode an jedem Punkt, während Grid Search nur eine stellt?
  Welche mathematische Größe liefert die Antwort auf die zweite Frage? \hfill (2~P)
  \item Grid Search benötigt $O(N^d)$ Auswertungen für $N$ Gitterpunkte in $d$ Dimensionen.
  Wie viele Auswertungen sind das für $N = 100$ Gitterpunkte pro Achse in $d = 6$
  Dimensionen? Wie heißt dieses Phänomen? \hfill (1~P)
  \item Unter welchem Namen ist \glqq Newton's method simplified to first order\grqq{} im
  Machine Learning bekannt, und warum ist diese Idee dort so zentral? \hfill (1~P)
\end{enumerate}

\schreibraum[7]
\Aufgabe{Herleitung der Newton-Raphson-Formel aus der linearen Approximation}{5}

Ziel der Nullstellensuche (root finding): Finde $x^*$ mit $f(x^*) = 0$. Ausgangspunkt ist
die lineare Approximation von $f$ nahe einem Punkt $x_n$ (Gl.~(20) im Skript):
\[ f(x) \approx f(x_n) + f'(x_n)\cdot(x - x_n) \qquad \text{für } x \text{ nahe } x_n. \]

\begin{enumerate}[label=(\alph*)]
  \item Warum ist es eine sinnvolle Strategie, die \emph{nächste} Schätzung $x_{n+1}$
  genau dort zu wählen, wo diese lineare Approximation die Null kreuzt? \hfill (1~P)
  \item Leiten Sie aus dem Ansatz
  \[ 0 = f(x_n) + f'(x_n)\cdot(x_{n+1} - x_n) \]
  die Newton-Raphson-Iterationsformel (Gl.~(21)) her. Geben Sie alle
  Umformungsschritte an. \hfill (2~P)
  \item Was bedeutet die Formel geometrisch? Beschreiben Sie die Rolle der Tangente
  (tangent line) am Punkt $(x_n, f(x_n))$. \hfill (1~P)
  \item Die Taylor-Entwicklung erlaubt beliebig hohe Approximationsordnungen. Warum
  begnügt sich Newtons Methode trotzdem mit der \emph{linearen} (1.~Ordnung)?
  Nennen Sie das Kosten-Nutzen-Argument des Skripts (\glqq Why linear is enough\grqq).
  \hfill (1~P)
\end{enumerate}

\schreibraum[8]
\Aufgabe{Newton-Iteration von Hand: $f(x) = x^2 - 5$}{7}

Berechnen Sie $\sqrt{5}$ mit Newtons Methode, d.\,h.\ finden Sie die positive Nullstelle
von $f(x) = x^2 - 5$. Starten Sie bei $x_0 = 3$.

\begin{enumerate}[label=(\alph*)]
  \item Geben Sie $f'(x)$ an und schreiben Sie die Newton-Iterationsvorschrift (Gl.~(21))
  konkret für diese Funktion hin. \hfill (1~P)
  \item Führen Sie drei Iterationsschritte durch ($x_1$, $x_2$, $x_3$). Tragen Sie
  Ihre Ergebnisse in eine Tabelle mit den Spalten $n$, $x_n$, $f(x_n)$, $f'(x_n)$,
  $x_{n+1}$ ein (sieben Nachkommastellen). \hfill (4~P)
  \item Berechnen Sie die Fehler $|e_n| = |x_n - \sqrt{5}|$ für $n = 0, \dots, 3$
  (Referenzwert $\sqrt{5} \approx 2.2360680$). Prüfen Sie zahlenmäßig, ob
  $|e_{n+1}| \approx C\,|e_n|^2$ mit einer ungefähr konstanten Konstanten $C$ gilt
  --- welcher Wert von $C$ ergibt sich? \hfill (2~P)
\end{enumerate}

\schreibraum[10]
\Aufgabe{Taylor-Entwicklung Ordnung für Ordnung: $f(x) = \ln(x)$ um $x_n = 1$}{6}

Das Skript baut die Taylor-Entwicklung (Taylor expansion) schrittweise auf: 0.~Ordnung
matcht den Funktionswert, 1.~Ordnung zusätzlich die erste Ableitung, 2.~Ordnung zusätzlich
die zweite Ableitung.

\begin{enumerate}[label=(\alph*)]
  \item Stellen Sie für $f(x) = \ln(x)$ die Approximationspolynome $P_0(x)$, $P_1(x)$
  und $P_2(x)$ um den Entwicklungspunkt $x_n = 1$ auf (Gl.~(23)--(25)). Geben Sie dazu
  zunächst $f(1)$, $f'(1)$ und $f''(1)$ an. \hfill (3~P)
  \item Werten Sie alle drei Polynome an der Stelle $x = 1.2$ aus und vergleichen Sie
  mit dem exakten Wert $\ln(1.2) \approx 0.1823216$: Geben Sie die drei absoluten Fehler
  an. Was beobachten Sie? \hfill (2~P)
  \item Das Skript zieht eine Analogie zwischen der Approximation 0.~Ordnung und einem
  Suchverfahren aus Kapitel~3. Welches Verfahren ist gemeint, und warum passt die
  Analogie? \hfill (1~P)
\end{enumerate}

\schreibraum[9]
\Aufgabe{Sekantenverfahren (secant method): Konzept und Handrechnung}{7}

Das Sekantenverfahren ersetzt die Ableitung in Newtons Formel durch den
Differenzenquotienten aus den zwei jüngsten Punkten (\glqq a poor man's derivative\grqq)
und liefert die Iteration (Gl.~(34)):
\[ x_{n+1} = x_n - f(x_n)\cdot\frac{x_n - x_{n-1}}{f(x_n) - f(x_{n-1})}. \]

\begin{enumerate}[label=(\alph*)]
  \item Warum benötigt das Sekantenverfahren \emph{zwei} Startwerte $x_0$ und $x_1$,
  während Newton mit einem auskommt? Argumentieren Sie über die Formel und über die
  geometrische Bedeutung der Sekante (secant line). \hfill (1~P)
  \item In welchen drei Situationen würden Sie das Sekantenverfahren gegenüber Newton
  bevorzugen (die drei Gründe des Skripts)? Und warum würden Sie es \emph{nicht}
  bevorzugen, wenn die Ableitung billig und exakt verfügbar ist? \hfill (2~P)
  \item Wenden Sie das Sekantenverfahren auf $f(x) = x^2 - 5$ mit den Startwerten
  $x_0 = 2$ und $x_1 = 3$ an. Berechnen Sie $x_2$, $x_3$ und $x_4$ mit allen
  Zwischenwerten ($f$-Werte, Differenzenquotient bzw.\ Schrittweite; sieben
  Nachkommastellen). Wie groß ist der Fehler $|x_4 - \sqrt{5}|$? \hfill (4~P)
\end{enumerate}

\schreibraum[10]
\Aufgabe{Quadratische Konvergenz quantitativ}{5}

Newtons Methode konvergiert nahe einer einfachen Nullstelle quadratisch:
$|e_{n+1}| \leq C\,|e_n|^2$ (Gl.~(33)).

\begin{enumerate}[label=(\alph*)]
  \item \emph{(Selbstreflexionsfrage des Skripts)} Wie viele Iterationen benötigt
  Newtons Methode, um von einem Fehler $10^{-1}$ auf einen Fehler $10^{-16}$ zu kommen?
  Nehmen Sie zur Vereinfachung $C \approx 1$ an und geben Sie die Fehlerfolge explizit
  an. Warum nennt man dieses Verhalten auch \glqq Verdopplung der gültigen Stellen pro
  Iteration\grqq? \hfill (2~P)
  \item Das Skript berechnet $\sqrt{2}$ mit Newton ab $x_0 = 2$; die Fehlerfolge
  lautet
  \begin{gather*}
    |e_0| = 5.86 \times 10^{-1},\qquad |e_1| = 8.58 \times 10^{-2},\qquad
    |e_2| = 2.45 \times 10^{-3},\\
    |e_3| = 2.12 \times 10^{-6},\qquad |e_4| \approx 1.6 \times 10^{-12}.
  \end{gather*}
  [Korrektur ggü.\ Original, S.~42: das PDF druckt in der letzten Zeile
  $4.5 \times 10^{-12}$ --- das ist das Residuum $|x_4^2 - 2|$, nicht der Fehler
  $|x_4 - \sqrt{2}|$; außerdem $|e_3| = 2.13 \times 10^{-6}$ statt
  $2.12 \times 10^{-6}$.]
  Prüfen Sie für jeden Schritt, ob $|e_{n+1}|$ ungefähr die Größenordnung von
  $|e_n|^2$ hat. Gegen welchen Wert stabilisiert sich das Verhältnis
  $|e_{n+1}| / |e_n|^2$? Vergleichen Sie mit der theoretischen Konstante $C$ aus
  Gl.~(32) (vgl.\ Aufgaben~\ref{auf:3} und~\ref{auf:8}).
  \hfill (2~P)
  \item Eine Randnotiz des Skripts zeigt die Fehlerfolge
  $10^{-1} \to 10^{-2} \to 10^{-3} \to 10^{-6} \to 10^{-12}$. Warum ist das Verhalten
  in den ersten beiden Schritten \emph{nicht} quadratisch, danach aber schon? Welche
  Annahme der Konvergenz-Herleitung steckt dahinter? \hfill (1~P)
\end{enumerate}

\schreibraum[8]
\Aufgabe{Transfer --- Geometrical Failure Search an $f(x) = x^3 - x$}{9}

\emph{(Angelehnt an die Original-Übung \glqq Geometrical Failure Search\grqq, S.~45.)}
Die Funktion $f(x) = x^3 - x$ hat die drei Nullstellen $x = -1, 0, 1$; ihre Ableitung ist
$f'(x) = 3x^2 - 1$. Zielwurzel sei $x = 1$. Finden Sie Startwerte $x_0$ für die drei
Fehlermodi (failure modes) von Newtons Methode --- und begründen Sie sie analytisch.

\begin{enumerate}[label=(\alph*)]
  \item \textbf{Nullableitung (zero derivative):} Bestimmen Sie \emph{alle} Startwerte
  $x_0$ mit $f'(x_0) = 0$. Was passiert dort geometrisch mit der Tangente, und warum
  bricht Algorithmus~1 des Skripts in diesem Fall ab? \hfill (2~P)
  \item \textbf{Konvergenz zur falschen Wurzel:} Zeigen Sie durch Rechnung, dass der
  Startwert $x_0 = 0.5$ --- der von den beiden Wurzeln $0$ und $1$ jeweils genau $0.5$
  entfernt ist --- in einem einzigen Newton-Schritt \emph{exakt} auf der \emph{entferntesten}
  Wurzel $x = -1$ landet, also weder die Zielwurzel $1$ noch die nächstgelegene
  Wurzelumgebung erreicht. \hfill (3~P)
  \item \textbf{Oszillation:} Die Newton-Abbildung lautet hier
  $N(x) = x - \frac{x^3 - x}{3x^2 - 1}$. Vereinfachen Sie $N(x)$ zu einem einzigen
  Bruch und bestimmen Sie daraus alle Startwerte $x_0 \neq 0$ mit $N(x_0) = -x_0$,
  bei denen Newton also für immer zwischen $x_0$ und $-x_0$ hin- und herspringt.
  Verifizieren Sie Ihren Wert durch Einsetzen für einen Iterationsschritt. \hfill (4~P)
\end{enumerate}

\schreibraum[12]
\Aufgabe{Transfer --- Herleitung der quadratischen Konvergenz skizzieren}{7}

Leiten Sie die Fehlerrekursion $|e_{n+1}| \leq C\,|e_n|^2$ entlang der Skript-Herleitung
(Gl.~(28)--(33)) her. Sei $x^*$ die Wurzel ($f(x^*) = 0$) und $e_n = x_n - x^*$ der
Fehler bei Iteration $n$.

\begin{enumerate}[label=(\alph*)]
  \item Schreiben Sie die Taylor-Entwicklung 2.~Ordnung von $f$ um die \emph{wahre
  Wurzel} $x^*$, ausgewertet an der Stelle $x_n$, mit Lagrange-Restglied
  (Zwischenstelle $\xi$ zwischen $x_n$ und $x^*$) auf, und vereinfachen Sie mit
  $f(x^*) = 0$ zu einem Ausdruck für $f(x_n)$ in $e_n$. \hfill (2~P)
  \item Zeigen Sie zunächst $e_{n+1} = e_n - \frac{f(x_n)}{f'(x_n)}$ und setzen Sie
  dann Ihren Ausdruck aus (a) ein. \hfill (2~P)
  \item Nutzen Sie die Annahme $f'(x_n) \approx f'(x^*)$ (Punkte nahe der Wurzel),
  um zu zeigen, dass sich der Fehlerterm \emph{erster} Ordnung aufhebt und
  $e_{n+1} \approx -\frac{f''(\xi)}{2 f'(x^*)}\, e_n^2$ übrig bleibt. Geben Sie an,
  wovon die Konstante $C$ in $|e_{n+1}| \leq C\,|e_n|^2$ abhängt. \hfill (2~P)
  \item Welche zwei Voraussetzungen müssen erfüllt sein, damit diese Herleitung
  trägt? (Stichworte: Art der Wurzel, Lage des Startwerts.) \hfill (1~P)
\end{enumerate}

\schreibraum[10]
\newpage

\section*{Lösungen --- erst nach Bearbeitung lesen!}

\subsection*{Lösung zu Aufgabe~\ref{auf:1}}

\begin{enumerate}[label=(\alph*)]
  \item Grid Search erkundet den Lösungsraum \emph{blind}: An jedem Gitterpunkt fragt es
  nur \glqq Was ist hier der Wert?\grqq{} Newtons Methode fragt zusätzlich
  \glqq In welche Richtung soll ich als Nächstes gehen?\grqq{} Die Antwort auf die zweite
  Frage liefert die \textbf{Ableitung} $f'(x_n)$ (die Steigung der Tangente): Sie zeigt
  Richtung \emph{und} sinnvolle Schrittweite zur Lösung. Der Hauptvorteil ist also die
  Nutzung \emph{lokaler Information} für gezielte Schritte statt blinder Auswertung
  --- und damit deutlich schnellere Konvergenz (im Skript-Beispiel: Newton erreicht mit
  4 Funktions-/Ableitungsauswertungen einen Fehler $< 10^{-5}$, Grid Search mit 6
  Auswertungen nur $0.042$).
  \item $N^d = 100^6 = 10^{12}$, also eine Billion Auswertungen. Das Phänomen heißt
  \emph{Fluch der Dimensionalität} (curse of dimensionality): Der Aufwand wächst
  exponentiell mit der Dimension $d$.
  \item Als \textbf{Gradientenabstieg} (gradient descent). Das gesamte Feld Deep Learning
  baut auf der Idee auf, mit lokaler Gradienteninformation in hochdimensionalen
  Loss-Landschaften zu Minima zu navigieren --- so werden Modelle auf spezifische
  Zielfunktionen hin optimiert (Backpropagation verteilt dabei die
  Gradienteninformation durch die Schichten des Netzes).
\end{enumerate}

\subsection*{Lösung zu Aufgabe~\ref{auf:2}}

\begin{enumerate}[label=(\alph*)]
  \item Nahe $x_n$ ist die lineare Approximation (Tangente) die beste verfügbare lokale
  Beschreibung von $f$. Wo die Tangente die Null kreuzt, ist daher \glqq our best guess
  for where the actual function crosses zero\grqq{} --- die beste Schätzung dafür, wo
  die \emph{tatsächliche} Funktion die Null kreuzt. Man löst also das (triviale) lineare
  Ersatzproblem exakt, statt das schwierige nichtlineare Problem direkt.
  \item Ansatz: Die lineare Approximation soll an der Stelle $x_{n+1}$ null sein:
  \begin{align*}
    0 &= f(x_n) + f'(x_n)\cdot(x_{n+1} - x_n) && \big| -f(x_n)\\
    f'(x_n)\cdot(x_{n+1} - x_n) &= -f(x_n) && \big| \, : f'(x_n) \ (\neq 0)\\
    x_{n+1} - x_n &= -\frac{f(x_n)}{f'(x_n)} && \big| + x_n\\[2pt]
    x_{n+1} &= x_n - \frac{f(x_n)}{f'(x_n)}. && \text{(Gl.~(21))}
  \end{align*}
  \emph{Hinweis [Korrektur ggü.\ Original, S.~38]:} Im PDF ist Gl.~(20) durch einen
  Layoutfehler verstümmelt (\glqq$+\,f(x_n)$\grqq{} hinter die Nebenbedingung gerutscht);
  die hier verwendete Form $f(x) \approx f(x_n) + f'(x_n)(x - x_n)$ ist die korrekte und
  wird durch die Folgezeile im Skript bestätigt.
  \item Geometrisch: Man legt die \textbf{Tangente} an den Graphen im Punkt
  $(x_n, f(x_n))$; ihr Schnittpunkt mit der $x$-Achse ist die nächste Schätzung
  $x_{n+1}$. Newton ersetzt also in jedem Schritt die Kurve durch ihre Tangente und
  springt zu deren Nullstelle.
  \item \glqq We want to solve $f(x^*) = 0$, not compute $f(x)$ as best as we can.\grqq{}
  Höhere Ordnungen kosten: höhere Ableitungen berechnen und kompliziertere Polynome
  auswerten (deren Nullstellen zudem schwerer zu finden sind). In numerischen Rechnungen
  muss die Taylor-Reihe ohnehin bei endlicher Ordnung abgeschnitten werden --- und der
  \emph{lineare} Term ist der erste nicht-konstante Term, der Richtungsinformation
  liefert. Er genügt, um pro Schritt eine gute nächste Schätzung zu erhalten.
\end{enumerate}

\subsection*{Lösung zu Aufgabe~\ref{auf:3}}

\begin{enumerate}[label=(\alph*)]
  \item $f'(x) = 2x$, also
  \[ x_{n+1} = x_n - \frac{x_n^2 - 5}{2 x_n}. \]
  \item Schrittweise (alle Werte auf sieben Nachkommastellen gerundet):
  \begin{align*}
    x_1 &= 3 - \frac{3^2 - 5}{2 \cdot 3} = 3 - \frac{4}{6} = 3 - 0.6666667 = 2.3333333,\\[2pt]
    x_2 &= 2.3333333 - \frac{2.3333333^2 - 5}{2 \cdot 2.3333333}
         = 2.3333333 - \frac{0.4444444}{4.6666667} = 2.2380952,\\[2pt]
    x_3 &= 2.2380952 - \frac{2.2380952^2 - 5}{2 \cdot 2.2380952}
         = 2.2380952 - \frac{0.0090703}{4.4761905} = 2.2360689.
  \end{align*}
  \begin{center}
  \begin{tabular}{ccccc}
    \toprule
    $n$ & $x_n$ & $f(x_n)$ & $f'(x_n)$ & $x_{n+1}$\\
    \midrule
    0 & 3.0000000 & 4.0000000 & 6.0000000 & 2.3333333\\
    1 & 2.3333333 & 0.4444444 & 4.6666667 & 2.2380952\\
    2 & 2.2380952 & 0.0090703 & 4.4761905 & 2.2360689\\
    3 & 2.2360689 & 0.0000041 & 4.4721378 & ---\\
    \bottomrule
  \end{tabular}
  \end{center}
  Bereits $x_3 = 2.2360689$ stimmt mit $\sqrt{5} = 2.2360680$ auf sechs
  Nachkommastellen überein.
  \item Fehler $|e_n| = |x_n - \sqrt{5}|$:
  \[ |e_0| = 7.639 \times 10^{-1}, \quad |e_1| = 9.727 \times 10^{-2}, \quad
     |e_2| = 2.027 \times 10^{-3}, \quad |e_3| = 9.18 \times 10^{-7}. \]
  Quotiententest $C_n = |e_{n+1}| / |e_n|^2$:
  \begin{gather*}
    C_0 = \frac{9.727 \times 10^{-2}}{(7.639 \times 10^{-1})^2} \approx 0.167, \qquad
    C_1 = \frac{2.027 \times 10^{-3}}{(9.727 \times 10^{-2})^2} \approx 0.214,\\
    C_2 = \frac{9.18 \times 10^{-7}}{(2.027 \times 10^{-3})^2} \approx 0.223.
  \end{gather*}
  Die Quotienten stabilisieren sich bei $C \approx 0.22$. Das passt exakt zur Theorie
  (Gl.~(32)): $C = \left|\frac{f''(x^*)}{2 f'(x^*)}\right| = \frac{2}{2 \cdot 2\sqrt{5}}
  = \frac{1}{2\sqrt{5}} \approx 0.2236$ --- quadratische Konvergenz bestätigt.
\end{enumerate}

\subsection*{Lösung zu Aufgabe~\ref{auf:4}}

\begin{enumerate}[label=(\alph*)]
  \item Ableitungen: $f(x) = \ln(x)$, $f'(x) = \frac{1}{x}$, $f''(x) = -\frac{1}{x^2}$,
  also $f(1) = 0$, $f'(1) = 1$, $f''(1) = -1$. Damit:
  \begin{align*}
    P_0(x) &= f(1) = 0 && \text{(0.~Ordnung, Gl.~(23): Funktionswert matchen)}\\
    P_1(x) &= f(1) + f'(1)(x - 1) = x - 1 && \text{(1.~Ordnung, Gl.~(24): Steigung matchen)}\\
    P_2(x) &= (x - 1) - \tfrac{1}{2}(x - 1)^2 && \text{(2.~Ordnung, Gl.~(25): Krümmung matchen)}
  \end{align*}
  (Der Faktor $\frac{1}{2}$ kompensiert die $2$, die beim zweimaligen Ableiten von
  $(x - x_n)^2$ entsteht --- Randnotiz \glqq Why divide by 2?\grqq.)
  \item Auswertung bei $x = 1.2$, also $(x-1) = 0.2$:
  \begin{align*}
    P_0(1.2) &= 0, & |P_0 - \ln(1.2)| &\approx 0.1823216,\\
    P_1(1.2) &= 0.2, & |P_1 - \ln(1.2)| &\approx 0.0176784,\\
    P_2(1.2) &= 0.2 - \tfrac{1}{2}(0.04) = 0.18, & |P_2 - \ln(1.2)| &\approx 0.0023216.
  \end{align*}
  Beobachtung: Jede zusätzliche Ordnung verkleinert den Fehler etwa um eine
  Größenordnung (Faktor $\approx 10$ bzw.\ $\approx 7.6$) --- die Taylor-Entwicklung ist
  eine \emph{lokale} Approximation, die nahe am Entwicklungspunkt mit jeder Ordnung
  rasch besser wird.
  \item Gemeint ist \textbf{Grid Search}: Die Approximation 0.~Ordnung nutzt nur den
  Funktionswert am Punkt selbst --- \glqq it only looks at the function value at
  discrete points, without any information about how the function behaves between those
  points.\grqq{} Genau so wertet Grid Search blind Punkt für Punkt aus, ohne
  Steigungsinformation.
\end{enumerate}

\subsection*{Lösung zu Aufgabe~\ref{auf:5}}

\begin{enumerate}[label=(\alph*)]
  \item Der Differenzenquotient
  $f'(x_n) \approx \frac{f(x_n) - f(x_{n-1})}{x_n - x_{n-1}}$ braucht \emph{zwei}
  verschiedene Stützstellen --- mit nur einem Punkt lässt sich keine Steigung schätzen.
  Geometrisch: Eine Sekante ist die Gerade durch \emph{zwei} Kurvenpunkte
  $(x_{n-1}, f(x_{n-1}))$ und $(x_n, f(x_n))$; erst ab dem zweiten Schritt liefert
  jede Iteration den jeweils zweiten Punkt automatisch mit.
  \item Sekante bevorzugen, wenn (die drei Skript-Gründe):
  \begin{itemize}[nosep]
    \item die Ableitung \emph{teuer} zu berechnen ist (the derivative is expensive to
    compute),
    \item die Funktion als \emph{Black Box} vorliegt, z.\,B.\ eine Simulation (the
    function is given as a black box),
    \item die Funktion nicht überall \emph{differenzierbar} ist (the function isn't
    differentiable everywhere).
  \end{itemize}
  Ist $f'$ dagegen billig und exakt verfügbar, ist Newton vorzuziehen: Die Sekante ist
  nur eine Näherung der Tangente, das Verfahren konvergiert daher etwas langsamer als
  quadratisch und benötigt typischerweise mehr Iterationen.
  \item Mit $f(x_0) = f(2) = -1$ und $f(x_1) = f(3) = 4$:
  \begin{align*}
    x_2 &= 3 - 4 \cdot \frac{3 - 2}{4 - (-1)} = 3 - \frac{4}{5} = 2.2,
        & f(x_2) &= 2.2^2 - 5 = -0.16,\\[2pt]
    x_3 &= 2.2 - (-0.16) \cdot \frac{2.2 - 3}{-0.16 - 4} = 2.2307692,
        & f(x_3) &= -0.0236686.
  \end{align*}
  Mit $x_3 - x_2 = 0.0307692$ und $f(x_3) - f(x_2) = -0.0236686 + 0.16 = 0.1363314$:
  \[ x_4 = 2.2307692 - (-0.0236686) \cdot \frac{0.0307692}{0.1363314}
         = 2.2307692 + 0.0053419 = 2.2361111. \]
  Fehler: $|x_4 - \sqrt{5}| = |2.2361111 - 2.2360680| \approx 4.3 \times 10^{-5}$.
  Das Verfahren konvergiert --- erkennbar langsamer als Newton (Aufgabe~\ref{auf:3}
  erreicht nach drei Schritten bereits $9.2 \times 10^{-7}$), dafür ganz ohne
  Ableitungsauswertung.

  \emph{Querbezug:} Das Skript rechnet dieselbe Methode an $f(x) = x^3 - x$ mit
  $x_0 = 1.2$, $x_1 = 1.4$ vor. Korrekte Werte dort: $f(1.2) = 0.528$,
  $f(1.4) = 1.344$ und $x_2 = 1.4 - 1.344 \cdot \frac{0.2}{0.816} \approx 1.0706$
  [Korrektur ggü.\ Original, S.~45: das PDF druckt $f(1.2) = 0.128$, $f(1.4) = 0.744$
  und darauf aufbauend eine falsche Iterationskette].
\end{enumerate}

\subsection*{Lösung zu Aufgabe~\ref{auf:6}}

\begin{enumerate}[label=(\alph*)]
  \item Mit $C \approx 1$ quadriert sich der Fehler in jedem Schritt --- der
  Zehnerexponent \emph{verdoppelt} sich:
  \[ 10^{-1} \;\to\; 10^{-2} \;\to\; 10^{-4} \;\to\; 10^{-8} \;\to\; 10^{-16}. \]
  Es genügen also \textbf{4 Iterationen}. Da der Fehler $10^{-k}$ ungefähr $k$ gültigen
  Dezimalstellen entspricht, verdoppelt jede Iteration die Zahl der gültigen Stellen
  ($1 \to 2 \to 4 \to 8 \to 16$) --- deshalb ist quadratische Konvergenz so mächtig:
  Maschinengenauigkeit (double precision, $\approx 10^{-16}$) ist in einer Handvoll
  Schritten erreicht.
  \item Schrittweise Prüfung von $|e_n|^2$ gegen $|e_{n+1}|$:
  \begin{center}
  \begin{tabular}{cccc}
    \toprule
    $n$ & $|e_n|^2$ & $|e_{n+1}|$ & Verhältnis $|e_{n+1}| / |e_n|^2$\\
    \midrule
    0 & $3.43 \times 10^{-1}$ & $8.58 \times 10^{-2}$ & $0.250$\\
    1 & $7.36 \times 10^{-3}$ & $2.45 \times 10^{-3}$ & $0.333$\\
    2 & $6.00 \times 10^{-6}$ & $2.12 \times 10^{-6}$ & $0.353$\\
    3 & $4.49 \times 10^{-12}$ & $\approx 1.6 \times 10^{-12}$ & $\approx 0.35$\\
    \bottomrule
  \end{tabular}
  \end{center}
  In jedem Schritt hat $|e_{n+1}|$ dieselbe Größenordnung wie $|e_n|^2$ ---
  quadratische Konvergenz von Anfang an erkennbar (Tabellenunterschrift im Skript:
  \glqq See how the error approximately squares each iteration.\grqq). Das Verhältnis
  $|e_{n+1}| / |e_n|^2$ stabilisiert sich dabei nicht bei irgendeinem Wert, sondern
  genau bei der theoretischen asymptotischen Konstante aus Gl.~(32): mit ungerundeten
  Werten lautet die Folge $0.250 \to 0.333 \to 0.353 \to 0.354$, und
  \[ C = \left|\frac{f''(\sqrt{2})}{2 f'(\sqrt{2})}\right|
     = \frac{2}{2 \cdot 2\sqrt{2}} = \frac{1}{2\sqrt{2}} \approx 0.3536 \]
  (Herleitung in Aufgabe~\ref{auf:8}; völlig analog zur Konstante
  $1/(2\sqrt{5}) \approx 0.2236$ beim $\sqrt{5}$-Beispiel in
  Aufgabe~\ref{auf:3}).
  \item In den ersten Schritten ($10^{-1} \to 10^{-2} \to 10^{-3}$) verdoppelt sich der
  Exponent \emph{nicht} --- der Fehler sinkt nur etwa linear um eine Größenordnung.
  Grund: Die Herleitung der quadratischen Konvergenz benutzt die Annahme
  $f'(x_n) \approx f'(x^*)$, die nur \emph{nahe genug an der Wurzel} gilt. Erst wenn
  die Iterierten in diesem Einzugsbereich sind (hier ab Iteration~2), greift die
  Quadrierung: $10^{-3} \to 10^{-6} \to 10^{-12}$ (Randnotiz: \glqq Once we're close
  enough and our assumptions hold, the error squares perfectly.\grqq).
\end{enumerate}

\subsection*{Lösung zu Aufgabe~\ref{auf:7}}

\begin{enumerate}[label=(\alph*)]
  \item $f'(x_0) = 3x_0^2 - 1 = 0 \iff x_0^2 = \frac{1}{3} \iff
  x_0 = \pm\frac{1}{\sqrt{3}} \approx \pm 0.5774$.
  An diesen Punkten ist die Tangente \emph{horizontal} --- sie kreuzt die $x$-Achse nie,
  der Newton-Schritt $-f(x_0)/f'(x_0)$ wäre eine Division durch null. Algorithmus~1
  fängt genau das ab: Bei $|d| < \epsilon_{\text{machine}}$ bricht er mit
  \glqq Derivative too small\grqq{} ab (sonst Division durch null bzw.\ numerische
  Instabilität).
  \item Mit $x_0 = 0.5$:
  \[ f(0.5) = 0.5^3 - 0.5 = 0.125 - 0.5 = -0.375, \qquad
     f'(0.5) = 3 \cdot 0.25 - 1 = -0.25, \]
  \[ x_1 = 0.5 - \frac{-0.375}{-0.25} = 0.5 - 1.5 = -1. \]
  Wegen $f(-1) = (-1)^3 - (-1) = 0$ landet die Iteration \emph{exakt} auf der Wurzel
  $x = -1$ und bleibt dort stehen. Obwohl $x_0 = 0.5$ genau mittig zwischen den Wurzeln
  $0$ und $1$ liegt, springt Newton zur entferntesten Wurzel $-1$: Die flache, negative
  Steigung ($f'(0.5) = -0.25$ nahe der Nullableitungsstelle $1/\sqrt{3} \approx 0.577$)
  erzeugt einen sehr großen Schritt. Das illustriert den Fehlermodus
  \glqq mehrdeutige Startpunktwahl\grqq: $x_0$ beeinflusst stark, welche Wurzel gefunden
  wird --- und es ist keineswegs immer die nächstgelegene.
  \item Vereinfachen auf einen Bruch:
  \[ N(x) = x - \frac{x^3 - x}{3x^2 - 1}
          = \frac{x(3x^2 - 1) - (x^3 - x)}{3x^2 - 1}
          = \frac{3x^3 - x - x^3 + x}{3x^2 - 1}
          = \frac{2x^3}{3x^2 - 1}. \]
  Oszillationsbedingung $N(x_0) = -x_0$ (mit $x_0 \neq 0$):
  \[ \frac{2x_0^3}{3x_0^2 - 1} = -x_0
     \iff 2x_0^3 = -x_0(3x_0^2 - 1) = -3x_0^3 + x_0
     \iff 5x_0^3 = x_0
     \iff x_0^2 = \frac{1}{5}, \]
  also $x_0 = \pm\frac{1}{\sqrt{5}} \approx \pm 0.4472$.
  Verifikation für $x_0 = \frac{1}{\sqrt{5}}$:
  \[ N\!\left(\tfrac{1}{\sqrt{5}}\right)
     = \frac{2 \cdot \frac{1}{5\sqrt{5}}}{\frac{3}{5} - 1}
     = \frac{\frac{2}{5\sqrt{5}}}{-\frac{2}{5}}
     = -\frac{1}{\sqrt{5}} \approx -0.4472136. \]
  Aus Symmetrie ($f$ ist ungerade, also $N(-x) = -N(x)$) folgt
  $N(-\tfrac{1}{\sqrt{5}}) = +\tfrac{1}{\sqrt{5}}$: Newton zykelt für immer
  $0.4472136 \to -0.4472136 \to 0.4472136 \to \cdots$ --- der Fehlermodus
  \emph{Oszillation}, der besonders bei symmetrischen Funktionen auftritt.

  \emph{Zum Vergleich (gutartiger Fall):} Vom Skript-Startwert $x_0 = 1.4$ konvergiert
  Newton dagegen monoton von oben gegen die Zielwurzel $x = 1$:
  $1.4 \to 1.1246 \to 1.0181 \to 1.0005$ [Korrektur ggü.\ Original, S.~44: das PDF
  druckt die fehlerhafte Kette $1.127 \to 0.976 \to 1.014$, die fälschlich unter die
  Wurzel springt].
\end{enumerate}

\subsection*{Lösung zu Aufgabe~\ref{auf:8}}

\begin{enumerate}[label=(\alph*)]
  \item Taylor 2.~Ordnung um $x^*$, ausgewertet bei $x_n$, mit Lagrange-Restglied
  (Gl.~(28)):
  \[ f(x_n) = f(x^*) + f'(x^*)(x_n - x^*) + \frac{f''(\xi)}{2}(x_n - x^*)^2, \]
  wobei $\xi$ zwischen $x_n$ und $x^*$ liegt. Mit $f(x^*) = 0$ und $e_n = x_n - x^*$
  folgt (Gl.~(29)):
  \[ f(x_n) = f'(x^*)\, e_n + \frac{f''(\xi)}{2}\, e_n^2. \]
  \item Aus der Newton-Formel (Gl.~(21)):
  \[ e_{n+1} = x_{n+1} - x^* = x_n - \frac{f(x_n)}{f'(x_n)} - x^*
     = (x_n - x^*) - \frac{f(x_n)}{f'(x_n)} = e_n - \frac{f(x_n)}{f'(x_n)}. \]
  Einsetzen von (a) liefert (Gl.~(30)):
  \[ e_{n+1} = e_n - \frac{f'(x^*)\, e_n + \frac{f''(\xi)}{2}\, e_n^2}{f'(x_n)}. \]
  \item Für Punkte nahe $x^*$ gilt $f'(x_n) \approx f'(x^*)$ (Gl.~(31)):
  \[ e_{n+1} \approx e_n - \frac{f'(x^*)\, e_n + \frac{f''(\xi)}{2}\, e_n^2}{f'(x^*)}
     = e_n - e_n - \frac{f''(\xi)}{2 f'(x^*)}\, e_n^2
     = -\frac{f''(\xi)}{2 f'(x^*)}\, e_n^2. \]
  Der Fehlerterm \emph{erster} Ordnung hebt sich also exakt auf --- übrig bleibt der
  quadratische Term (Gl.~(32)), abstrakter $|e_{n+1}| \leq C\,|e_n|^2$ (Gl.~(33)) mit
  \[ C = \left|\frac{f''(\xi)}{2 f'(x^*)}\right|: \]
  $C$ hängt von der \emph{Krümmung} ($f''$ nahe der Wurzel) und der \emph{Steigung}
  ($f'$ an der Wurzel) ab. Große Krümmung oder flache Steigung vergrößern $C$ und
  verlangsamen die Konvergenz.
  \item Voraussetzungen:
  \begin{itemize}[nosep]
    \item \textbf{Einfache Wurzel} (simple root): $f'(x^*) \neq 0$ --- sonst Division
    durch null in $C$; bei mehrfachen Wurzeln degradiert Newton auf lineare Konvergenz.
    \item \textbf{Startwert nah genug} an $x^*$: Die Annahme $f'(x_n) \approx f'(x^*)$
    (und die Gültigkeit der lokalen Taylor-Approximation) greift nur in einer Umgebung
    der Wurzel. Außerhalb drohen die Fehlermodi aus Aufgabe~\ref{auf:7}
    (\glqq Newton-Raphson gives (limited) convergence guarantees near simple
    roots.\grqq).
  \end{itemize}
\end{enumerate}

\end{document}
